

% ---------------------------------------------------------------
\section{Enunciado}
% ---------------------------------------------------------------
Simule un flujo en el canal para distintos números de Reynolds (considere $\Rey = 100, 1000, 10000$ y $100000$) en la geometría indicada usando \OF{} (use condiciones \emph{inlet-outlet}). Haga una simulación de flujo incompresible transiente, y use un modelo de turbulencia cuando sea necesario.

% ---------------------------------------------------------------
\section{Geometría y datos del problema}
% ---------------------------------------------------------------
\subsection{Geometría}

\begin{figure}[H]
\centering
\begin{tikzpicture}[>=Stealth, scale=1]
  % Esquema fuera de escala
  \def\L{12}\def\H{3}
  % paredes superior e inferior
  \fill[gray!30] (0,\H) rectangle (\L,\H+0.25);
  \fill[gray!30] (0,-0.25) rectangle (\L,0);
  \draw[thick] (0,0) -- (\L,0);
  \draw[thick] (0,\H) -- (\L,\H);
  % fluido
  \fill[cyan!10] (0,0) rectangle (\L,\H);
  \draw[thick] (0,0) -- (\L,0);
  \draw[thick] (0,\H) -- (\L,\H);
  % entrada y salida
  \draw[dashed] (0,0) -- (0,\H);
  \draw[dashed] (\L,0) -- (\L,\H);
  % cilindro
  \filldraw[fill=gray!35,draw=black,thick] (2,{\H/2}) circle (0.4);
  % cotas
  \draw[<->] (\L+0.5,0) -- (\L+0.5,\H) node[midway,right] {$H = \SI{20}{m}$};
  \draw[<->] (0,-0.7) -- (\L,-0.7) node[midway,below] {$L = \SI{100}{m}$};
  \draw[<->] (0,{\H/2}) -- (1.6,{\H/2}) node[midway,above] {$\SI{5}{m}$};
  \draw[<->] (2.7,{\H/2-0.4}) -- (2.7,{\H/2+0.4}) node[right] {$D = \SI{2}{m}$};
  % flujo
  \draw[->,thick] (-1.6,{\H/2}) -- (-0.7,{\H/2}) node[midway,above] {$U$};
  \node at (0.35,{\H/2+0.7}) {\footnotesize inlet};
  \node at (\L-0.5,{\H/2+0.7}) {\footnotesize outlet};
\end{tikzpicture}
\caption{Esquema del dominio (fuera de escala), según el enunciado.}
\label{fig:geometria}
\end{figure}

\begin{table}[H]
\centering
\begin{tabular}{llc}
\toprule
Parámetro & Símbolo & Valor \\
\midrule
Largo del canal & $L$ & \SI{100}{m} \\
Alto del canal & $H$ & \SI{20}{m} \\
Diámetro del cilindro & $D$ & \SI{2}{m} \\
Distancia entrada--cilindro & $x_c$ & \SI{5}{m} \\
Razón de bloqueo & $D/H$ & $0{,}1$ \\
\bottomrule
\end{tabular}
\caption{Datos geométricos del problema.}
\end{table}

\paragraph{Supuestos adoptados.}
\begin{itemize}[nosep]
  \item La cota de \SI{5}{m} se mide desde la entrada hasta el \emph{borde} frontal del cilindro, por lo que el centro queda en $x = \SI{6}{m}$ (en \archivo{Parametros}: \texttt{Lup = 5 + R}). El resto del canal, hasta completar $L=\SI{100}{m}$, queda aguas abajo.
  \item El cilindro está centrado verticalmente ($y = \SI{10}{m}$).
  \item Las franjas grises superior e inferior se interpretan como paredes con condición de no deslizamiento.
  \item Problema 2D: una sola celda en la dirección $z$ con patch \texttt{empty}.
\end{itemize}

\subsection{Números de Reynolds y propiedades}
Se define $\Rey = U D/\nu$ con $D = \SI{2}{m}$. Fijando $U = \SI{1}{m/s}$ y variando la viscosidad cinemática $\nu$:

\begin{table}[H]
\centering
\begin{tabular}{rcc}
\toprule
$\Rey$ & $U$ [\si{m/s}] & $\nu = UD/\Rey$ [\si{m^2/s}] \\
\midrule
$10^{2}$ & 1 & $2\times10^{-2}$ \\
$10^{3}$ & 1 & $2\times10^{-3}$ \\
$10^{4}$ & 1 & $2\times10^{-4}$ \\
$10^{5}$ & 1 & $2\times10^{-5}$ \\
\bottomrule
\end{tabular}
\caption{Parámetros de cada caso (con $U$ fijo, se varía $\nu$ en \texttt{transportProperties}).}
\end{table}

% ---------------------------------------------------------------
\section{Plan de trabajo}
% ---------------------------------------------------------------
Lista de todo lo que hay que hacer para resolver el problema:

\begin{enumerate}[label=\textbf{\arabic*.}]
  \item \textbf{Definir los casos}: un directorio por Reynolds (\texttt{Re100}, \texttt{Re1000}, \texttt{Re10000}, \texttt{Re100000}) (\texttt{Cilindro\_Re<N>}); $\Rey=10^2$ y $10^3$ comparten la misma malla.
  \item \textbf{Malla 2D} (\texttt{blockMesh} + \texttt{snappyHexMesh} + \texttt{extrudeMesh}): malla de fondo del canal, refinamiento y capas prismáticas en el cilindro y caja de refinamiento en la estela (Sección~\ref{sec:refinamiento}). Patches: \texttt{inlet}, \texttt{outlet}, \texttt{lowerWall}, \texttt{upperWall}, \texttt{cylinder}, \texttt{front}/\texttt{back} (\texttt{empty}).
  \item \textbf{Calidad de malla}: \texttt{checkMesh}; estimar $y^+$ para elegir tratamiento de pared según el modelo de turbulencia.
  \item \textbf{Condiciones de borde} (ver Tabla~\ref{tab:bc}).
  \item \textbf{Modelo físico}: \texttt{pimpleFoam} (incompresible, transiente). Elegir turbulencia según $\Rey$ (ver Sección~\ref{sec:modelo}).
  \item \textbf{Paso de tiempo}: controlar Courant $\mathrm{Co} < 1$ (\texttt{adjustTimeStep}); tiempo total suficiente para varios periodos de desprendimiento de vórtices.
  \item \textbf{Post-proceso}: $\Cd$, $\Cl$ (\texttt{forceCoeffs}), $\St$ (FFT de $\Cl$), campos de velocidad/presión/vorticidad en ParaView, y comparación entre Reynolds.
  \item \textbf{Verificación}: sensibilidad a malla y a $\Delta t$; comparación con valores de la literatura.
  \item \textbf{Redactar el informe} con este documento.
\end{enumerate}

% ---------------------------------------------------------------
\section{Modelo y configuración numérica}
\label{sec:modelo}
% ---------------------------------------------------------------
\subsection{Ecuaciones}
Flujo incompresible, transiente, 2D:
\begin{equation}
  \nabla\cdot\mathbf{u} = 0, \qquad
  \frac{\partial \mathbf{u}}{\partial t} + \nabla\cdot(\mathbf{u}\mathbf{u})
  = -\frac{1}{\rho}\nabla p + \nabla\cdot\!\left[(\nu+\nu_t)\left(\nabla\mathbf{u}+\nabla\mathbf{u}^{T}\right)\right]
\end{equation}
donde $\nu_t = 0$ en los casos laminares.

\subsection{Modelo de turbulencia por caso}
\begin{table}[H]
\centering
\begin{tabular}{rll}
\toprule
$\Rey$ & Régimen esperado & Modelo \\
\midrule
$10^{2}$ & Laminar, estela periódica (sobre $\Rey\approx47$) & \texttt{laminar} \\
$10^{3}$ & Laminar 2D (en 3D real ya es inestable) & \texttt{laminar} \\
$10^{4}$ & Capa límite laminar, estela turbulenta & URANS 2D, $k$--$\omega$ SST \\
$10^{5}$ & Capa límite laminar, estela turbulenta & URANS 2D, $k$--$\omega$ SST \\
\bottomrule
\end{tabular}
\caption{Selección del modelo de turbulencia.}
\end{table}

\paragraph{Justificación.}
\begin{itemize}[nosep]
  \item \textbf{$\Rey = 10^2$.} Está sobre el $\Rey$ crítico de desprendimiento ($\approx 47$) pero bajo la aparición de inestabilidades 3D ($\approx 190$) \cite{williamson1996}. La simulación 2D laminar es físicamente consistente y resuelve la calle de vórtices de von Kármán sin modelo.
  \item \textbf{$\Rey = 10^3$.} Se mantiene laminar porque la malla, con celdas de \SI{0.025}{m} en la pared ($D/80$), es suficientemente fina para resolver las escalas. Es una simulación 2D de un flujo que en la realidad es tridimensional, por lo que cabe esperar un $\Cd$ y un $\Cl$ mayores que los experimentales (la tridimensionalidad reduce la coherencia de la estela).
  \item \textbf{$\Rey = 10^4, 10^5$.} Se usa $k$--$\omega$ SST \cite{menter1994}: se comporta como $k$--$\omega$ cerca de la pared (válido hasta la subcapa viscosa con $\yplus\approx 1$) y como $k$--$\varepsilon$ en el flujo libre, y maneja bien los gradientes adversos de presión y la separación. Es \emph{URANS}: separa la escala del desprendimiento (resuelta) de la turbulencia (modelada).
  \item \textbf{Limitación.} SST trata la capa límite como turbulenta desde el inicio, mientras que en estos $\Rey$ (subcríticos, bajo la crisis de arrastre en $\Rey\approx 2$--$3\times10^{5}$) la capa límite real es laminar y separa antes ($\approx 80^\circ$). El modelo tenderá a retrasar la separación y a subestimar $\Cd$. Las condiciones de entrada y de pared están en la Sección~\ref{sec:bc}.
\end{itemize}

\subsection{Condiciones de borde}
\label{sec:bc}
\begin{table}[H]
\centering
\footnotesize
\setlength{\tabcolsep}{4pt}
\begin{tabular}{lllll}
\toprule
Patch & $\mathbf{u}$ & $p$ & $k,\ \omega$ & $\nu_t$ \\
\midrule
\texttt{inlet} & \texttt{fixedValue} $(U,0,0)$ & \texttt{zeroGradient} & \texttt{fixedValue} & \texttt{calculated} \\
\texttt{outlet} & \texttt{inletOutlet} & \texttt{fixedValue} 0 & \texttt{inletOutlet} & \texttt{calculated} \\
\texttt{cylinder} & \texttt{noSlip} & \texttt{zeroGradient} & funciones de pared & \texttt{nutkWallFunction} \\
\texttt{lowerWall}, \texttt{upperWall} & \texttt{noSlip} & \texttt{zeroGradient} & funciones de pared & \texttt{nutkWallFunction} \\
\texttt{front}, \texttt{back} & \texttt{empty} & \texttt{empty} & \texttt{empty} & \texttt{empty} \\
\bottomrule
\end{tabular}
\caption{Condiciones de borde. Para $k$ se usa \texttt{kqRWallFunction} y para $\omega$ \texttt{omegaWallFunction}; los casos laminares no tienen $k$, $\omega$ ni $\nu_t$.}
\label{tab:bc}
\end{table}

Con $\yplus\approx 1$ en el cilindro, \texttt{omegaWallFunction} y \texttt{nutkWallFunction} se comportan como una condición de pared de bajo $\Rey$; en las paredes del canal, que no tienen capas, actúan como funciones de pared estándar.

\paragraph{Turbulencia en la entrada.} Se impone intensidad $I = \num{0.01}$ y razón de viscosidades $\nu_t/\nu = 10$:
\begin{equation}
  k_{\mathrm{in}} = \tfrac{3}{2}(U I)^2 = \SI{1.5e-4}{m^2/s^2}, \qquad
  \omega_{\mathrm{in}} = \frac{k_{\mathrm{in}}}{\nu_t} = \frac{k_{\mathrm{in}}}{10\,\nu}.
\end{equation}
Esto da $\omega_{\mathrm{in}} = \SI{0.075}{s^{-1}}$ para $\Rey = 10^4$ y $\SI{0.75}{s^{-1}}$ para $\Rey = 10^5$. La salida usa \texttt{inletOutlet} para evitar que un eventual reflujo introduzca energía.

\subsection{Malla}
\label{sec:malla}
La malla se genera en tres pasos con las utilidades de \OF{}, todo controlado desde un único archivo de parámetros (\archivo{constant/Parametros}) que leen todos los diccionarios:

\begin{enumerate}[nosep]
  \item \comando{blockMesh}: malla de fondo del canal completo, \emph{sin} cilindro, con celdas cúbicas de lado $\Delta x = D/10 = \SI{0.2}{m}$ ($500\times100$ celdas, $\num{50000}$ en total) y una sola celda en $z$.
  \item \comando{snappyHexMesh}: recorta el cilindro, refina alrededor y en la estela, ajusta (\emph{snap}) la malla a la superficie y agrega capas prismáticas en la pared del cilindro.
  \item \comando{extrudeMesh}: \comando{snappyHexMesh} también refina en $z$, por lo que se toma la cara \texttt{front} y se extruye \emph{una} celda para recuperar una malla 2D con \texttt{front}/\texttt{back} de tipo \texttt{empty}.
\end{enumerate}

\subsubsection{Refinamiento cerca del cilindro}
\label{sec:refinamiento}
\paragraph{Por qué se necesita.} Los gradientes más fuertes del problema están pegados al cilindro: la capa límite, los puntos de separación y las capas de corte que luego se enrollan en los vórtices. Una malla uniforme de \SI{0.2}{m} ($D/10$) no los resolvería, pero refinar todo el canal ($\num{500}\times\num{100}$ celdas a $\Delta x$) con el tamaño necesario en la pared sería inviable. La solución es refinar solo donde se necesita, en tres zonas anidadas (\cref{fig:zonas}):

\begin{enumerate}[label=\textbf{\arabic*.}]
  \item \textbf{Superficie del cilindro} (refinamiento más fuerte).
  \item \textbf{Capas prismáticas} pegadas a la pared del cilindro.
  \item \textbf{Caja de estela}, con refinamiento moderado, donde se desprenden y se convectan los vórtices.
\end{enumerate}
Fuera de ellas la malla queda en el tamaño de fondo, y en las paredes del canal no se refina (allí se usan funciones de pared).

\paragraph{Niveles de refinamiento (octree).} \comando{snappyHexMesh} refina dividiendo cada celda en 8 (en 2D, en 4): un \emph{nivel} $n$ da celdas de tamaño
\begin{equation}
  \Delta_n = \frac{\Delta x}{2^{n}} = \frac{D}{10\cdot 2^{n}}.
\end{equation}
\begin{table}[H]
\centering
\begin{tabular}{cccl}
\toprule
Nivel $n$ & $\Delta_n$ [\si{m}] & Celdas por $D$ & Uso \\
\midrule
0 & \num{0.2}     & 10  & malla de fondo \\
2 & \num{0.05}    & 40  & estela, $\Rey\le10^3$ \\
3 & \num{0.025}   & 80  & estela, $\Rey\ge10^4$; cilindro, $\Rey\le10^3$ \\
4 & \num{0.0125}  & 160 & cilindro, $\Rey=10^4$ \\
5 & \num{0.00625} & 320 & cilindro, $\Rey=10^5$ \\
\bottomrule
\end{tabular}
\caption{Tamaños de celda por nivel de refinamiento.}
\end{table}

En \archivo{snappyHexMeshDict} esto se declara con
\begin{foamcode}
refinementSurfaces
{
    cylinder { level ($cylLevelMin $cylLevelMax); }   // en la superficie
}
refinementRegions
{
    wake { mode inside; level $wakeLevel; }           // dentro de la caja
}
nCellsBetweenLevels 3;     // 3 celdas de transicion entre niveles
\end{foamcode}
Se usa $\texttt{cylLevelMin} = \texttt{cylLevelMax}$ porque el cilindro es liso y no tiene curvatura variable que justifique un rango. \texttt{nCellsBetweenLevels = 3} fuerza una transición gradual: entre una zona de nivel $n$ y otra de nivel $n-1$ hay tres celdas de relleno, de modo que el tamaño cambia de a un factor 2 por vez y la razón de aspecto y la oblicuidad se mantienen bajas.

\paragraph{Adaptación a la superficie (\emph{snap}).} El refinamiento por octree deja la pared como una escalera de celdas. El paso \emph{snap} (\texttt{snapControls}) desplaza los nodos de esa escalera sobre el cilindro geométrico, que se declara como un cilindro analítico de radio $R = D/2$ (no como una malla de triángulos), por lo que la superficie queda exacta y sin facetas.

\paragraph{Capas prismáticas.} Con un tamaño de celda $\Delta_n$ aislado, no se puede alcanzar $\yplus\approx 1$ sin una malla enorme, porque la primera celda tendría que ser mucho más delgada que ancha. Por eso se agregan $N$ capas prismáticas (\texttt{addLayersControls}) pegadas a la pared, de espesor creciente geométricamente con razón $e_r$:
\begin{equation}
  t_i = t_1\, e_r^{\,i-1},\qquad
  T = \sum_{i=1}^{N} t_i = t_1\,\frac{e_r^{N}-1}{e_r-1},\qquad
  t_N = t_1\, e_r^{\,N-1}.
\end{equation}
Los espesores se dan en valores absolutos (\texttt{relativeSizes false}, \texttt{firstLayerThickness} $=t_1$), con $e_r = 1.2$ en todos los casos. Las capas se agregan solo en el cilindro; \texttt{featureAngle 180} porque la superficie es suave, y \texttt{minThickness} $=0.1\,t_1$ como espesor mínimo admisible.

\paragraph{Cálculo del primer espesor $t_1$.} Se impone que la primera capa tenga el espesor de $\yplus=1$. Con la velocidad de fricción $u_\tau = U\sqrt{C_f/2}$ se tiene $y_1 = \nu/u_\tau$. El coeficiente de fricción se estima en la mitad de la circunferencia ($x = \pi D/2$, $\Rey_x = \Rey\,\pi/2$) con correlaciones de placa plana: Blasius, $C_f = 1.328/\sqrt{\Rey_x}$, en los casos laminares, y $C_f = 0.026\,\Rey_x^{-1/7}$ en los turbulentos. Es una estimación de orden de magnitud, no un cálculo sobre el cilindro real; $\yplus$ se verificará con el resultado de la simulación. \pend{reportar $\yplus$ máximo y promedio sobre el cilindro desde \comando{postProcess -func yPlus}}

Para $\Rey = 10^2$ y $10^3$ el $y_1$ resultante es mayor que la propia celda ($y_1=\SI{0.087}{m}$ y $\SI{0.0155}{m}$ frente a $\Delta_3=\SI{0.025}{m}$), es decir, la capa límite laminar se resuelve sin necesidad de refinamiento especial, y se toma $t_1 = \Delta_n/2$.

\paragraph{Cálculo del número de capas $N$.} Se elige $N$ para que la última capa tenga aproximadamente el tamaño de la celda del octree vecina, $t_N\approx\Delta_n$. Así el paso de la zona prismática a la zona hexaédrica es suave y no hay un salto brusco de volumen. Despejando,
\begin{equation}
  N \approx 1 + \frac{\ln(\Delta_n/t_1)}{\ln e_r}.
\end{equation}

\begin{table}[H]
\centering
\small
\begin{tabular}{rlccccccc}
\toprule
$\Rey$ & Modelo & $n$ & $\Delta_n$ [\si{m}] & $N$ & $t_1$ [\si{m}] & $t_N$ [\si{m}] & $T$ [\si{m}] & $y_1$ [\si{m}] \\
\midrule
$10^{2}$ & laminar & 3 & \num{0.025}   & 5  & \num{0.0125}  & \num{0.0259} & \num{0.093}  & \num{0.087}  \\
$10^{3}$ & laminar & 3 & \num{0.025}   & 5  & \num{0.0125}  & \num{0.0259} & \num{0.093}  & \num{0.0155} \\
$10^{4}$ & SST     & 4 & \num{0.0125}  & 8  & \num{0.0035}  & \num{0.0125} & \num{0.058}  & \num{0.0035} \\
$10^{5}$ & SST     & 5 & \num{0.00625} & 16 & \num{0.00041} & \num{0.0063} & \num{0.036}  & \num{0.00041} \\
\bottomrule
\end{tabular}
\caption{Refinamiento en el cilindro. $y_1$: espesor para $\yplus=1$ (estimación de placa plana); $t_1$: espesor impuesto. Con $t_1\approx y_1$ el centro de la primera celda queda en $\yplus\approx0.5$.}
\label{tab:capas}
\end{table}

El espesor total $T$ es menor que el espesor de capa límite estimado (de orden \SIrange{0.1}{0.2}{m} para los casos turbulentos con correlaciones de placa plana), de modo que las capas quedan dentro de la capa límite. Para $\Rey=10^{2}$ la capa límite es comparable a $D$ y las capas cubren solo una fracción pequeña, lo que no es problema porque allí la malla del octree ya la resuelve.

\paragraph{Caja de estela.} El desprendimiento ocurre aguas abajo, así que se refina una caja
\[
  x\in[x_c-1.5D,\ x_c+8D],\qquad y\in[y_c-1.5D,\ y_c+1.5D],
\]
es decir $x\in[3,\,22]\,\si{m}$ e $y\in[7,\,13]\,\si{m}$, con nivel 2 ($\Rey\le10^3$) o 3 ($\Rey\ge10^4$). Su extensión cubre los primeros $\approx 8D$ de estela, donde se forma la calle de vórtices; más lejos la solución es suave y no justifica el costo. El nivel es menor que el de la pared porque la escala de los vórtices ya desprendidos es del orden de $D$, no de la capa límite.

\begin{figure}[H]
\centering
\begin{tikzpicture}[scale=0.12,>=Stealth]
  % dominio a escala: 100 x 20 m
  \draw[thick,fill=cyan!8] (0,0) rectangle (100,20);
  % caja de estela x in [3,22], y in [7,13]
  \draw[ofazul,thick,fill=ofazul!15] (3,7) rectangle (22,13);
  % cilindro
  \filldraw[fill=gray!50,draw=black,thick] (6,10) circle (1);
  % anillo de capas (exagerado)
  \draw[ofrojo,thick] (6,10) circle (1.7);
  % cotas
  \draw[<->] (0,-2.5) -- (100,-2.5) node[midway,below] {\footnotesize $L=\SI{100}{m}$};
  \draw[<->] (102.5,0) -- (102.5,20) node[midway,right] {\footnotesize $H=\SI{20}{m}$};
  \node[ofazul,font=\footnotesize,anchor=west] at (24,10) {caja de estela};
  \node[ofrojo,font=\footnotesize,anchor=south] at (6,13.7) {capas};
  \node[font=\footnotesize] at (60,17) {malla de fondo ($\Delta x=\SI{0.2}{m}$)};
  \draw[->,thick] (-9,10) -- (-4,10) node[midway,above] {\footnotesize $U$};
\end{tikzpicture}
\caption{Zonas de refinamiento (a escala; el espesor de las capas está exagerado para que se vea). La caja de estela cubre $x\in[3,22]\,\si{m}$.}
\label{fig:zonas}
\end{figure}

\begin{figure}[H]
\centering
\begin{tikzpicture}[scale=1.1]
  % Esquema (no a escala) de las capas prismaticas, N = 6, er = 1.2
  \fill[gray!35] (0,-0.3) rectangle (9,0);
  \draw[thick] (0,0) -- (9,0);
  \node[anchor=east] at (-0.1,-0.15) {\footnotesize pared (cilindro)};
  % alturas acumuladas: t1 = 0.5, er = 1.2
  \foreach \y in {0.5,1.1,1.82,2.684,3.721,4.965} {
    \draw[ofrojo] (0,\y) -- (9,\y);
  }
  % columnas de ancho ~ t_N
  \foreach \x in {1.244,2.488,...,8.7} {
    \draw[gray] (\x,0) -- (\x,4.965);
  }
  % celdas del octree por encima
  \foreach \x in {0,1.244,2.488,...,8.7} {
    \draw[ofazul] (\x,4.965) rectangle ++(1.244,1.244);
  }
  \draw[<->] (9.3,0) -- (9.3,0.5) node[midway,right] {\footnotesize $t_1$};
  \draw[<->] (9.3,3.721) -- (9.3,4.965) node[midway,right] {\footnotesize $t_N\approx\Delta_n$};
  \draw[<->] (9.9,0) -- (9.9,4.965) node[midway,right] {\footnotesize $T$};
  \node[ofazul,anchor=west,font=\footnotesize] at (0.1,6.6) {celdas del octree, $\Delta_n$};
\end{tikzpicture}
\caption{Esquema de las capas prismáticas ($N=6$, $e_r=1.2$): la primera capa es delgada ($\yplus\approx1$) y la última iguala el tamaño $\Delta_n$ de las celdas hexaédricas vecinas.}
\label{fig:capas}
\end{figure}

\subsubsection{Malla obtenida y calidad}
Los resultados de \comando{snappyHexMesh} y \comando{checkMesh} (en 2D, tras \comando{extrudeMesh}) son:
\begin{table}[H]
\centering
\begin{tabular}{rrcc}
\toprule
$\Rey$ & Celdas & Razón de aspecto máx. & Oblicuidad máx. \\
\midrule
$10^{2}$ & \num{94104}  & \num{2.35}  & \num{0.39} \\
$10^{3}$ & \num{94104}  & \num{2.35}  & \num{0.39} \\
$10^{4}$ & \num{232590} & \num{3.50}  & \num{0.41} \\
$10^{5}$ & \num{246474} & \num{14.5}  & \num{0.41} \\
\bottomrule
\end{tabular}
\caption{Tamaño y calidad de las mallas finales. \comando{checkMesh} reporta \texttt{Mesh OK} en los cuatro casos.}
\end{table}
Las capas se completaron al \SI{100}{\percent} del espesor objetivo en todos los casos (\texttt{log.snappyHexMesh}). La razón de aspecto de $\approx14.5$ para $\Rey=10^5$ es esperada: la primera capa mide $\SI{0.41}{mm}$ de alto y $\Delta_5=\SI{6.25}{mm}$ de ancho, es decir una celda alargada en la dirección tangencial, que es la práctica habitual en una capa límite (el gradiente importante es normal a la pared). Los casos $\Rey=10^2$ y $10^3$ comparten la misma malla.

\begin{figure}[H]
\centering
% \includegraphics[width=0.9\textwidth]{Imagenes/T2/malla_global.png}
\fbox{\parbox{0.8\textwidth}{\centering\vspace{1.5cm}\pend{captura de la malla completa en ParaView}\vspace{1.5cm}}}
\caption{Malla utilizada (global).}
\end{figure}

\begin{figure}[H]
\centering
% \includegraphics[width=0.9\textwidth]{Imagenes/T2/malla_cilindro.png}
\fbox{\parbox{0.8\textwidth}{\centering\vspace{1.5cm}\pend{captura del detalle de capas alrededor del cilindro}\vspace{1.5cm}}}
\caption{Detalle de la malla alrededor del cilindro, con las capas prismáticas.}
\end{figure}

\subsection{Configuración numérica}
Se usa \comando{foamRun} con el solver \texttt{incompressibleFluid} (\OF{} 14), que resuelve el sistema transiente con el algoritmo PIMPLE.

\begin{table}[H]
\centering
\small
\begin{tabular}{ll}
\toprule
Término & Esquema \\
\midrule
Temporal & \texttt{backward} (2.º orden) \\
Gradientes & \texttt{Gauss linear} \\
Convección de $\mathbf{u}$ & \texttt{Gauss linearUpwind grad(U)} (2.º orden) \\
Convección de $k$, $\omega$ & \texttt{Gauss limitedLinear 1} \\
Laplaciano / $\nabla_n$ & \texttt{Gauss linear corrected} / \texttt{corrected} \\
\midrule
$p$ & GAMG, $\texttt{relTol}=0.01$ (0 en el paso final) \\
$\mathbf{u}$, $k$, $\omega$ & \texttt{smoothSolver}, \texttt{symGaussSeidel} \\
PIMPLE & 1 corrector externo, 2 de presión, 1 no ortogonal \\
\bottomrule
\end{tabular}
\caption{Esquemas numéricos y solvers.}
\end{table}

\paragraph{Paso de tiempo y duración.} El paso se ajusta con $\Cou_{\max}=0.5$ (\texttt{adjustTimeStep}), partiendo de $\Delta t = 10^{-3}\,t_c$ con $t_c = D/U = \SI{2}{s}$. Cada caso corre $200\,t_c$ ($\SI{400}{s}$) y escribe cada $t_c$. Para el Strouhal se descartará el transitorio inicial.

\paragraph{Coeficientes de fuerza.} \texttt{forceCoeffs} sobre el patch \texttt{cylinder} con $\rho_\infty=1$, $U_\infty = \SI{1}{m/s}$, $l_{\mathrm{ref}} = D$ y $A_{\mathrm{ref}} = D\,\Delta z$, donde $\Delta z=\SI{0.2}{m}$ es el espesor de la malla en $z$; así $\Cd$ y $\Cl$ son por unidad de longitud adimensionalizados con $\tfrac12\rho U^2 D$. Además hay una sonda de $\mathbf{u}$ y $p$ a $2D$ aguas abajo del centro del cilindro.

\subsection{Archivos de configuración y reproducción}
Cada caso es un directorio \texttt{Cilindro\_Re<N>} con
\begin{lstlisting}
Cilindro_Re<N>/
  Allrun, Allclean
  0/            U, p, (k, omega, nut)
  constant/     Parametros, physicalProperties, momentumTransport
  system/       blockMeshDict, snappyHexMeshDict, extrudeMeshDict,
                controlDict, fvSchemes, fvSolution
\end{lstlisting}
El único archivo que cambia entre casos es \archivo{constant/Parametros}. Su bloque de refinamiento (para $\Rey=10^5$) es:
\begin{foamcode}
D       2;     H 20;    Lup #calc "5.0 + 0.5*$D";     // centro a 5 m + R del inlet
Re      100000;  Uin 1;  nu #calc "$Uin*$D/$Re";
dx      #calc "$D/10.0";                              // malla de fondo 0.2 m
wakeLevel 3;
cylLevelMin 5;  cylLevelMax 5;                        // dn = dx/2^5 = 0.00625 m
nLayers 16;  expansionRatio 1.2;  firstLayerThickness 0.00041;
\end{foamcode}
Para correr un caso (las simulaciones no forman parte de este documento):
\begin{terminal}
source /opt/openfoam14/etc/bashrc
cd Cilindro_Re<N>
./Allrun     # blockMesh, snappyHexMesh, extrudeMesh, checkMesh, foamRun
\end{terminal}


% ---------------------------------------------------------------
\section{Resultados}
% ---------------------------------------------------------------
Para cada $\Rey$ incluir: campo de velocidad, presión y vorticidad en un instante representativo, y la historia temporal de $\Cd$ y $\Cl$.

\subsection{\texorpdfstring{$\Rey = 100$}{Re = 100}}
\pend{figuras y comentarios}
\subsection{\texorpdfstring{$\Rey = 1000$}{Re = 1000}}
\pend{figuras y comentarios}
\subsection{\texorpdfstring{$\Rey = 10000$}{Re = 10000}}
\pend{figuras y comentarios}
\subsection{\texorpdfstring{$\Rey = 100000$}{Re = 100000}}
\pend{figuras y comentarios}

\subsection{Comparación entre casos}
\begin{table}[H]
\centering
\begin{tabular}{rccc}
\toprule
$\Rey$ & $\overline{\Cd}$ & $\Cl^{\mathrm{rms}}$ & $\St$ \\
\midrule
$10^{2}$ & \pend{} & \pend{} & \pend{} \\
$10^{3}$ & \pend{} & \pend{} & \pend{} \\
$10^{4}$ & \pend{} & \pend{} & \pend{} \\
$10^{5}$ & \pend{} & \pend{} & \pend{} \\
\bottomrule
\end{tabular}
\caption{Coeficientes integrales y número de Strouhal, con $\St = f D/U$.}
\end{table}

% ---------------------------------------------------------------
\section{Análisis y discusión}
% ---------------------------------------------------------------
\pend{efecto del Reynolds sobre la estela, desprendimiento de vórtices, influencia del bloqueo $D/H = 0{,}1$ y de las paredes, rol del modelo de turbulencia y limitaciones 2D/URANS, sensibilidad a malla y paso de tiempo}

% ---------------------------------------------------------------
\section{Conclusiones}
% ---------------------------------------------------------------
\pend{conclusiones}

\begin{thebibliography}{9}
\bibitem{menter1994} F.~R. Menter, ``Two-equation eddy-viscosity turbulence models for engineering applications'', \emph{AIAA Journal}, 32(8):1598--1605, 1994.
\bibitem{williamson1996} C.~H.~K. Williamson, ``Vortex dynamics in the cylinder wake'', \emph{Annual Review of Fluid Mechanics}, 28:477--539, 1996.
\bibitem{schlichting} H.~Schlichting y K.~Gersten, \emph{Boundary-Layer Theory}, Springer, 9.ª ed.
\bibitem{openfoam} OpenFOAM Foundation, \emph{OpenFOAM User Guide}, versión 14, \url{https://doc.cfd.direct/openfoam/user-guide-v14/}.
\end{thebibliography}

